\documentclass[a4paper,11pt]{article}
\usepackage[utf8]{inputenc}
\usepackage{geometry}
\usepackage{setspace}
\usepackage{lmodern}
\usepackage{graphicx}
\usepackage{array}
\usepackage{booktabs}
\usepackage{titlesec}
\usepackage{fancyhdr}
\usepackage{xcolor}
\usepackage{colortbl}
\usepackage[spanish]{babel}
\usepackage{hyperref}

% --- CONFIGURACIÓN DE PÁGINA ---
\geometry{left=2.5cm, right=2.5cm, top=2.5cm, bottom=2.5cm}
\setstretch{1.2}
\pagestyle{fancy}
\fancyhf{}
\fancyhead[L]{\textcolor{gray}{WASAFF CONSULTING}}
\fancyhead[R]{\textcolor{gray}{Propuesta: Diagnóstico Estratégico de Riesgos}}
\fancyfoot[C]{\thepage}

% --- COLORES Y ESTILOS ---
\definecolor{WasaffBlue}{RGB}{41, 52, 98}
\definecolor{WasaffGray}{RGB}{100, 100, 100}
\titleformat{\section}{\Large\bfseries\color{WasaffBlue}}{\thesection}{1em}{}
\linespread{1.3}

% ===================================================================================
\begin{document}

\begin{titlepage}
    \centering
    \vspace*{1cm}
    {\scshape\Huge\bfseries WASAFF CONSULTING \par}
    \vspace{0.2cm}
    {\large\itshape De la Física a la Solución.\par}
    \vspace{2cm}
    \rule{\linewidth}{0.8pt}
    \vspace{0.4cm}
    {\Huge\bfseries Diagnóstico Estratégico de Riesgos\par}
    \vspace{0.1cm}
    {\large\itshape Maniobra y Suspensión de Tuberías de Acero\par}
    \rule{\linewidth}{0.2pt}
    \vspace{2.5cm}
    
    \begin{tabular}{ll}
        \textbf{Preparado para:} & Sr. Francisco Pérez \\
                                & Synergy Fire Protection \\
        \textbf{Fecha:}           & 22 de agosto de 2024 \\
        \textbf{Propuesta Nº:}    & \textbf{DIAG-WC-2024-001} \\
        \textbf{Validez:}         & 15 días \\
    \end{tabular}
    
    \vfill
    {\small Santiago, Chile \textbar\ +56 9 4612 5682 \textbar\ \href{mailto:felipe@wasaffconsulting.com}{felipe@wasaffconsulting.com} \textbar\ \url{www.wasaffconsulting.com}}
\end{titlepage}

\section*{1. Entendimiento del Desafío}
Estimado Sr. Pérez,

Gracias por su confianza. Entendemos que el desafío de Synergy Fire Protection va más allá de la simple instalación de tuberías; se trata de \textbf{garantizar la seguridad de sus equipos y la continuidad operacional de sus proyectos}. Una falla en la maniobra de suspensión de tuberías de acero, especialmente en un rango amplio de diámetros y pesos, no es solo un incidente, sino un riesgo significativo que puede derivar en accidentes, detenciones de obra y sobrecostos imprevistos.

Nuestra propuesta no es un informe de ingeniería tradicional. Es un \textbf{Diagnóstico Estratégico de 10 días}, diseñado para entregarle una \textbf{hoja de ruta cuantificada}. Usaremos modelado y simulación basados en primeros principios para transformar la incertidumbre en una evaluación de riesgos clara, permitiéndole tomar decisiones con datos, no con conjeturas.

\section*{2. Nuestro Proceso: La Pregunta a la Respuesta Cuantificada}
En 10 días hábiles, seguiremos un proceso ágil y enfocado:

\begin{enumerate}
    \item \textbf{Sesión de Inmersión (Día 1-2):} Alinearemos el alcance preciso del diagnóstico, confirmando los diámetros (1'' a 12''), Schedules (10 a 40) y alturas de maniobra (1 a 16 m) que representan el mayor riesgo para su operación.
    
    \item \textbf{Análisis y Modelado (Día 3-7):} Construiremos un "laboratorio virtual" del escenario. Usando las herramientas de la mecánica fundamental, simularemos los escenarios de caída para cuantificar las energías de impacto y los radios de afectación. \textit{Aquí aplicamos el poder de los Capítulos 4 (Elasticidad) y la dinámica de cuerpos rígidos.}
    
    \item \textbf{Identificación y Cuantificación de Escenarios (Día 8-9):} Traduciremos los resultados complejos en 2 o 3 escenarios de riesgo priorizados. Para cada uno, cuantificaremos la probabilidad y el impacto potencial, estableciendo las bases para un análisis ROI de las posibles soluciones de mitigación.
    
    \item \textbf{Entrega de Hoja de Ruta (Día 10):} Le entregaremos y presentaremos un informe ejecutivo con los hallazgos.
\end{enumerate}

\section*{3. El Entregable: Claridad y Dirección}
Al final de los 10 días hábiles, usted recibirá:
\begin{itemize}
    \item Un análisis claro de los escenarios de mayor riesgo en sus maniobras.
    \item Una \textbf{cuantificación del impacto potencial} (energía liberada, radios de seguridad recomendados).
    \item \textbf{Recomendaciones técnicas priorizadas} para la mitigación (ej. protocolos de izaje, zonas de exclusión, EPP específico).
    \item Una base sólida y auditable para justificar la implementación de nuevos protocolos de seguridad ante sus clientes y equipos.
\end{itemize}
El informe incluirá gráficos, tablas en formato XLSX y el código de simulación utilizado.

\section*{4. Inversión: Valor Fijo y Transparente}
Creemos que el primer paso hacia una solución no debe ser una inversión impredecible. Por eso, nuestro Diagnóstico Estratégico se ofrece con un valor fijo, alineado con nuestra oferta de servicio.

\begin{table}[h!]
\centering
\begin{tabular}{lrr}
    \toprule
    \textbf{Concepto} & \textbf{Valor (UF)} & \textbf{Valor (CLP aprox.)} \\
    \midrule
    Diagnóstico Estratégico de Riesgos (Plazo 10 días hábiles) & 45.00 & \$1,695,514 \\
    \midrule
    \textbf{Subtotal} & \textbf{45.00} & \textbf{\$1,695,514} \\
    IVA (19\%) & 8.55 & \$322,148 \\
    \textbf{Total} & \textbf{53.55} & \textbf{\$2,017,662} \\
    \bottomrule
\end{tabular}
\caption{Valor de la inversión para la Hoja de Ruta. (1 UF = 37,678.09 CLP).}
\end{table}

\paragraph{Forma de Pago:} 50\% para iniciar el trabajo y 50\% contra entrega de la Hoja de Ruta final.

\paragraph{Nuestro Compromiso de Valor:} Entendemos que este diagnóstico es el primer paso. Si decide avanzar con nosotros a una fase de implementación o ingeniería de detalle, \textbf{el 50\% del valor de este diagnóstico se abonará como crédito a ese proyecto}. Su inversión inicial se transforma; nunca se pierde.

\vspace{1cm}
Agradecemos la oportunidad de presentarle esta propuesta. Estamos convencidos de que podemos entregarle la claridad y la dirección que su proyecto necesita.

\vspace{2cm}
\noindent
Atentamente,\\[2em]
\textbf{\textcolor{WasaffBlue}{Felipe Wasaff}}\\
Fundador, Wasaff Consulting\\
\href{mailto:felipe@wasaffconsulting.com}{felipe@wasaffconsulting.com}

\end{document}
